% P18-A1: finite abstract argumentation, serving N-x-009.
\chapter{Finite Abstract Argumentation}
\label{sec:p18-finite-argumentation}

An abstract argumentation framework records which arguments attack which
others. An extension selects arguments according to a stated rule for
conflict and defense. Different rules can give different selections, and
some rules can give no selection at all. This chapter develops the finite
foundations of Dung's framework and four extension semantics
\cite[secs.~2.1--2.2, pp.~325--330]{Dung1995}.

The inputs are a finite set of argument identifiers and a directed attack
relation. The results concern those inputs. They neither check an argument's
premises nor establish that its conclusion is true. Applying the framework
requires a separate account of how arguments and attacks were obtained.

\section{Attack and defense}
\label{sec:p18-attack-defense}

\begin{hbdefinition}{def:p18-finite-framework}{Finite argumentation framework}
A \emph{finite argumentation framework} is a pair
\(\mathcal F=(A,R)\), where \(A\) is a finite set of
\emph{arguments} and \(R\subseteq A\times A\) is its \emph{attack
relation}. The assertion \((a,b)\in R\) means that \(a\) attacks
\(b\). The set \(A\) may be empty. Attacks need not be symmetric,
and \((a,a)\in R\) is permitted: it is a \emph{self-attack}.

For \(a\in A\) and \(S\subseteq A\), define
\[
 \operatorname{Att}_{\mathcal F}(a)\defeq\{b\in A:(b,a)\in R\},
 \qquad
 S^+_{\mathcal F}\defeq\{b\in A:\exists s\in S\ (s,b)\in R\}.
\]
These are respectively the attackers of \(a\) and the arguments attacked
by \(S\). A set \(S\) is \emph{conflict-free} if there are no
\(a,b\in S\), including \(a=b\), with \((a,b)\in R\).
Thus the condition is \(S\cap S^+_{\mathcal F}=\varnothing\).
For a fixed framework the subscripts may be omitted.
\index{argumentation framework!finite}\index{attack relation}
\end{hbdefinition}

\begin{hbdefinition}{def:p18-defense-operator}{Defense and the characteristic operator}
A set \(S\subseteq A\) \emph{defends} an argument \(a\in A\)
when every attacker of \(a\) is attacked by some member of \(S\):
\[
 \operatorname{Att}_{\mathcal F}(a)\subseteq S^+_{\mathcal F}.
\]
Dung calls such an argument \emph{acceptable with respect to \(S\)}.
An argument with no attackers is defended by every set, including the empty
set. Defense of \(a\) does not require \(a\in S\).

The \emph{characteristic operator} is the function
\[
 \mathsf F_{\mathcal F}:2^A\longrightarrow2^A,\qquad
 \mathsf F_{\mathcal F}(S)
 \defeq\{a\in A:\operatorname{Att}_{\mathcal F}(a)\subseteq S^+_{\mathcal F}\},
\]
where \(2^A\) is the set of all subsets of \(A\). Write \(\mathsf F\)
when the framework is fixed. A set \(S\) is \emph{admissible} when it
is conflict-free and defends all its members, equivalently when it is
conflict-free and \(S\subseteq\mathsf F(S)\).
\index{argumentation!defense}\index{argumentation!admissible set}
\index{characteristic operator!argumentation}
\end{hbdefinition}

\begin{hbtheorem}{thm:p18-defense-properties}{proved=annex:p18.defense-properties}{Monotonicity and admissible enlargement}
The operator \(\mathsf F\) is monotone with respect to inclusion:
\(S\subseteq T\) implies \(\mathsf F(S)\subseteq\mathsf F(T)\).
If \(S\) is conflict-free, then \(\mathsf F(S)\) is conflict-free.
If \(S\) is admissible, then \(\mathsf F(S)\) is admissible.

In particular, if \(S\) is admissible and \(a\in\mathsf F(S)\),
then \(S\cup\{a\}\) is admissible. Every argument defended by \(S\)
remains defended by \(S\cup\{a\}\). This is the finite instance of
Dung's Fundamental Lemma, with self-attacks included
\cite[lemma~10 and lemmas~18--19, pp.~327--329]{Dung1995}.
\end{hbtheorem}

\section{Extension semantics}
\label{sec:p18-extension-semantics}

\begin{hbdefinition}{def:p18-extension-semantics}{Complete, grounded, preferred and stable extensions}
For a fixed \(\mathcal F=(A,R)\):
\begin{enumerate}
\item A \emph{complete extension} is an admissible set \(S\) containing every
      argument it defends: \(\mathsf F(S)\subseteq S\).
\item A \emph{fixed point} of \(\mathsf F\) is a set \(S\) with
      \(\mathsf F(S)=S\). Its \emph{least fixed point}, if it exists,
      is a fixed point contained in every fixed point. This least fixed
      point is the \emph{grounded extension} \(G_{\mathcal F}\).
\item A \emph{preferred extension} is an admissible set maximal under
      inclusion: it has no strictly larger admissible superset.
\item A \emph{stable extension} is a conflict-free set \(S\) that attacks
      every argument outside it: \(A\setminus S\subseteq S^+\).
\end{enumerate}
The families of complete, preferred and stable extensions are written
\(\operatorname{Comp}(\mathcal F)\),
\(\operatorname{Pref}(\mathcal F)\), and
\(\operatorname{Stab}(\mathcal F)\). The grounded-extension family is
\(\{G_{\mathcal F}\}\) once existence has been established.

Maximality of a preferred extension concerns containing supersets, not
maximum cardinality among all admissible sets. Leastness of the grounded
extension means containment in every relevant set; it is stronger than
merely having no smaller fixed point. The existence statements follow below.
\index{extension!complete}\index{extension!grounded}
\index{extension!preferred}\index{extension!stable}
\end{hbdefinition}

\begin{hbtheorem}{thm:p18-grounded-iteration}{proved=annex:p18.grounded-iteration}{Finite grounded iteration}
Put \(n=|A|\), \(G_0=\varnothing\), and
\(G_{i+1}=\mathsf F(G_i)\) for every integer \(i\geq0\).
Every \(G_i\) is admissible and \(G_i\subseteq G_{i+1}\).
There is \(k\leq n\) with \(G_k=G_{k+1}\).
This fixed point is the unique least fixed point of \(\mathsf F\), so
the grounded extension exists and equals \(G_k\).

A set is complete exactly when it is a conflict-free fixed point of
\(\mathsf F\). The grounded extension is complete and is contained in
every complete extension. Hence complete extensions always exist, and
\[
 G_{\mathcal F}=\bigcap_{E\in\operatorname{Comp}(\mathcal F)}E.
\]
\end{hbtheorem}

\begin{hbmethod}{def:p18-grounded-procedure}{Computing the grounded extension}
Given the finite argument list and attack relation, put \(n=|A|\) and
start with \(S=\varnothing\).
Form \(S^+\) by collecting all targets of attacks whose source is in
\(S\). Next form
\(T=\{a\in A:\operatorname{Att}(a)\subseteq S^+\}\).
If \(T=S\), return \(S\); otherwise replace \(S\) by \(T\)
and repeat. Theorem~\ref{thm:p18-grounded-iteration} proves correctness and
termination after at most \(n+1\) evaluations of \(\mathsf F\), counting
the final equality check. The bound is on operator evaluations, not a claim
that evaluating the operator takes constant time.
\end{hbmethod}

\begin{hbtheorem}{thm:p18-extension-relations}{proved=annex:p18.extension-relations}{Existence and inclusion relations}
Every admissible set in a finite framework is contained in a preferred
extension. In particular, a preferred extension always exists.
Every preferred extension is complete, and every stable extension is
preferred. Thus
\[
\begin{aligned}
 \operatorname{Stab}(\mathcal F)
 &\subseteq\operatorname{Pref}(\mathcal F)
  \subseteq\operatorname{Comp}(\mathcal F),\\
 G_{\mathcal F}&\subseteq E\qquad
   (E\in\operatorname{Comp}(\mathcal F)).
\end{aligned}
\]
A set \(S\subseteq A\) is stable exactly when
\(S=A\setminus S^+\). A stable extension need not exist. Preferred
extensions need not be stable, and complete extensions need not be preferred;
the examples below witness these failures of converse inclusion.
\end{hbtheorem}

\begin{hbwarning}{def:p18-fixed-point-boundary}{A fixed point can contain a conflict}
The equation \(S=\mathsf F(S)\) alone is not the definition of a
complete extension: conflict-freeness is also required. In the one-argument
self-attack framework below, \(\{a\}\) is a fixed point but is not
conflict-free. The grounded construction starts from the admissible empty
set and preserves that condition at every stage.
\end{hbwarning}

\section{Small frameworks}
\label{sec:p18-small-frameworks}

\begin{hbdefinition}{def:p18-example-frameworks}{Six finite frameworks}
The following pairs explicitly list all arguments and all attacks; omitted
attacks are absent. The case names are labels for these mathematical inputs.
\[
\begin{aligned}
 \mathcal F_{\rm chain}
   &=(\{a,b,c\},\{(a,b),(b,c)\}),\\
 \mathcal F_{\rm pair}
   &=(\{a,b\},\{(a,b),(b,a)\}),\\
 \mathcal F_{\rm self}
   &=(\{a\},\{(a,a)\}),\\
 \mathcal F_{\rm empty}
   &=(\varnothing,\varnothing),\\
 \mathcal F_{\rm size}
   &=(\{a,b,c\},\{(a,b),(b,a),(a,c),(c,a)\}),\\
 \mathcal F_{\rm shared}
   &=(\{a,b,c,d\},\{(a,b),(b,a),(a,c),(b,c),(c,d)\}).
\end{aligned}
\]
The mutual pair has Dung's Nixon-diamond attack pattern, while the self-attack
is his counterexample to the converse of stable-implies-preferred
\cite[example~9 and lemma~15, pp.~327--328]{Dung1995}.
The other named cases provide additional finite illustrations here.
\end{hbdefinition}

\begin{hbtheorem}{thm:p18-example-extensions}{proved=annex:p18.example-extensions}{Extension families in the six frameworks}
The complete lists of extensions for the defined cases are:
\begin{enumerate}
\item \emph{Chain.} The grounded stages are
      \(\varnothing,\{a\},\{a,c\},\{a,c\}\). Each of the three
      families \(\operatorname{Comp},\operatorname{Pref},\operatorname{Stab}\)
      consists only of \(\{a,c\}\).
\item \emph{Mutual pair.} The grounded extension is empty, and
      \[
      \begin{aligned}
       \operatorname{Comp}(\mathcal F_{\rm pair})
          &=\{\varnothing,\{a\},\{b\}\},\\
       \operatorname{Pref}(\mathcal F_{\rm pair})
          =\operatorname{Stab}(\mathcal F_{\rm pair})
          &=\{\{a\},\{b\}\}.
      \end{aligned}
      \]
\item \emph{Self-attack.} The grounded extension is empty,
      \(\operatorname{Comp}(\mathcal F_{\rm self})
        =\operatorname{Pref}(\mathcal F_{\rm self})=\{\varnothing\}\),
      and \(\operatorname{Stab}(\mathcal F_{\rm self})=\varnothing\).
      Nevertheless \(\mathsf F_{\mathcal F_{\rm self}}(\{a\})=\{a\}\).
\item \emph{Empty framework.} The grounded extension is empty and all
      three families equal \(\{\varnothing\}\). Thus there is one
      stable extension, whose argument set is empty.
\item \emph{Unequal sizes.} The grounded extension is empty, and
      \[
      \begin{aligned}
       \operatorname{Comp}(\mathcal F_{\rm size})
          &=\{\varnothing,\{a\},\{b,c\}\},\\
       \operatorname{Pref}(\mathcal F_{\rm size})
          =\operatorname{Stab}(\mathcal F_{\rm size})
          &=\{\{a\},\{b,c\}\}.
      \end{aligned}
      \]
      Both preferred extensions are maximal under inclusion, although their
      cardinalities differ.
\item \emph{Shared argument.} The grounded extension is empty,
      and
      \[
      \begin{aligned}
       \operatorname{Comp}(\mathcal F_{\rm shared})
          &=\{\varnothing,\{a,d\},\{b,d\}\},\\
       \operatorname{Pref}(\mathcal F_{\rm shared})
          =\operatorname{Stab}(\mathcal F_{\rm shared})
          &=\{\{a,d\},\{b,d\}\}.
      \end{aligned}
      \]
      The argument \(d\) belongs to every preferred extension but is absent
      from the grounded extension.
\end{enumerate}
\end{hbtheorem}

\section{Acceptance queries and limits}
\label{sec:p18-acceptance-boundary}

\begin{hbdefinition}{def:p18-acceptance-queries}{Credulous and skeptical acceptance queries}
Choose an extension family \(\mathcal E\subseteq2^A\), such as
\(\operatorname{Pref}(\mathcal F)\). An argument \(a\in A\) is
\emph{credulously accepted} when it belongs to at least one member of
\(\mathcal E\). In this chapter's query convention it is
\emph{skeptically accepted} when \(\mathcal E\ne\varnothing\) and
it belongs to every member of \(\mathcal E\).

The explicit nonemptiness guard is part of this query convention. When
\(\mathcal E=\varnothing\), report \emph{no extensions}; neither
acceptance label is assigned to an argument. An unguarded universal
quantifier over that family would instead be vacuously true for every
argument. When \(\mathcal E=\{\varnothing\}\), an extension exists
but contains no arguments. These two cases are distinct.
\index{acceptance!credulous}\index{acceptance!skeptical}
\end{hbdefinition}

The source calls grounded semantics skeptical and preferred semantics
credulous. Those names do not identify the grounded extension with the
intersection of all preferred extensions. The shared-argument case above
exhibits the difference noted by Dung
\cite[remark~26, p.~330]{Dung1995}. Acceptance must therefore name its chosen
extension family. The guarded query convention also does not modify the
definition of a stable extension or force one to exist.

\begin{hbwarning}{def:p18-framework-interpretation}{What an extension establishes}
An extension result establishes a property of the declared arguments and
attacks under one of the stated semantics. It does not establish the truth
of the arguments' premises, the correctness of a claimed attack, the
credibility of a speaker, or the empirical adequacy of the selected semantics.
In this framework an argument is an identifier, and no function assigning
premises or conclusions to identifiers has been supplied.

Infinite frameworks, finitary versus non-finitary behavior and the full
transfinite constructions are supplied in
\cref{sec:p18-arbitrary-argumentation}. That treatment also proves the
source's arbitrary-framework graph conditions and alternating-level
comparison. Structured arguments, dialogue protocols, further semantic
variants and the source's embeddings of other logics and games retain
their distinct treatments.
\end{hbwarning}
